\section{Operations budget and cost details}\label{sec:upbudget}

This replaces the USDF half of the operations budget (\tabref{tab:opsSummary},
\tabref{tab:opsSumUSDF}) and the USDF cost details of \secref{sec:opsdetails}.
The Chilean budget and cost details, and the cloud costs (\tabref{tab:cloud}),
are not updated.

Two forecasts are given, both priced with USDF hardware quotes:
\begin{itemize}
\item \emph{All at the USDF} (\tabref{tab:smCost}): every workload is hosted
      and bought at the USDF. It is an upper bound for the USDF.
\item \emph{Shared DRP} (\tabref{tab:smCostSplit}): the USDF's own purchases
      when DRP processing is shared with France and the UK at the shares in
      \secref{sec:upsplit}. It is the realistic USDF budget, and the one to
      compare with \tabref{tab:opsSumUSDF}, which likewise deducted the
      partners' share of DRP compute.
\end{itemize}
No price quotes are held for the partner facilities, so their hardware is
priced in neither forecast. Each partner buys its own.

\subsection{Method}

Both forecasts use the same purchase rules:

\begin{itemize}
\item \emph{Batch compute}, for DRP only: new Torino nodes when
      Milano-equivalent peak demand exceeds the fleet, plus cohort refresh.
\item \emph{Disk}: bought to the on-floor need, with re-purchase at the end
      of the storage lifetime.
\item \emph{Tape}: bought once the cumulative footprint exceeds the capacity
      assumed available.
\item \emph{Kubernetes}: a fixed annual growth rate, plus cohort refresh.
\item \emph{Qserv}: bought to the running maximum of the node projection,
      plus cohort refresh.
\end{itemize}

\smpara{What the shared-DRP forecast changes.}
Only two demands differ:
\begin{itemize}
\item Batch compute follows the USDF's share of peak cores.
\item On-floor disk excludes the live intermediates held at the partner
      facilities during their share of each campaign.
\end{itemize}
Everything else stays at the USDF, and those lines are identical in both
forecasts:
\begin{itemize}
\item Prompt Processing, the Data Access Centre and the other Kubernetes
      services;
\item the USDF's Qserv;
\item raw data and the retained data-release products, on disk and on tape.
\end{itemize}
Tape backup is a USDF responsibility in the multi-site DRP plan. Partner
facilities also run Qserv deployments of their own (DP2 has been ingested at
the UKDF, and a Polish Data Facility deployment is planned); those are partner
costs. The other and miscellaneous disk allowance is not reduced, because the
user areas, APDB and simulation output it stands in for stay at the USDF
whoever does the processing.

\smpara{Refresh by cohort.}
Hardware bought in fiscal year $x$ is replaced in year $x$ plus its lifetime,
starting from the purchases in \tabref{tab:smPreops}. The lifetime is
\smLifeCompute{} years for compute and \smLifeStorage{} for storage. The
construction-era tables refreshed compute after three years.

\smpara{FY2026 uses actual requests.}
The first column of both forecasts uses the FY2026 requests recorded in the
2026 capacity capture, not modelled figures.

\smpara{Unit costs} come from the 2026 capacity capture (August 2026
revision), and replace the Xeon and Rome exemplars of
\tabref{tab:opsXeonUSDF} and \tabref{tab:opsRomeUSDF}:
\begin{itemize}
\item Torino batch node with 768\,GB: \smTorinoUSD{}.
\item Kubernetes or Qserv-class node: \smKeightsUSD{}.
\item Usable Ceph disk: \smCephPerTB{} per TB.
\item Tape: \$\smTapePerTB{} per TB.
\end{itemize}
Prices are held flat. The construction-era model assumed annual falls of 10\%
for CPU, 5\% for storage and 8\% for Qserv (\tabref{tab:Inputs}).

There is no current price for the fast (NVMe) tier of \tabref{tab:Storage}, so
that tier is sized in \secref{sec:upstorage} but not separately costed. Hosting
and overhead rates, which replace the NCSA overheads of
\tabref{tab:opsOverheadCost}, were carried forward from an earlier agreement
rather than re-quoted. They are the least well-founded cost inputs.

\subsection{Result}

\begin{center}
\begin{tabular}{lrr}
\hline
Ten-year total & All at the USDF & Shared DRP (USDF only) \\
\hline
Batch compute & \smTenYearBatch{} & \smSplitTenYearBatch{} \\
Disk & \smTenYearDisk{} & \smSplitTenYearDisk{} \\
Tape & \smTenYearTape{} & \smTenYearTape{} \\
Kubernetes & \smTenYearKeights{} & \smTenYearKeights{} \\
Qserv & \smTenYearQserv{} & \smTenYearQserv{} \\
\hline
Hardware & \smTenYearHardware{} & \smSplitTenYearHardware{} \\
Total, with hosting and overhead & \smTenYearTotal{} & \smSplitTenYearTotal{} \\
\hline
\end{tabular}
\end{center}

Tape, Kubernetes and Qserv are the same in both forecasts, because they stay
at the USDF. The choice of DRP compute estimate is the largest single
sensitivity. With all processing at the USDF, moving from the superseded
\wfortyone{} estimate to \vthirty{} (\secref{sec:upcompute}) adds roughly
\$29M over ten years.

The forecast is expected to trigger the project risk on a computing cost
mismatch (Rubin-DP-T-9). The responses available include:
\begin{itemize}
\item accepting longer processing;
\item further efficiency in algorithms, storage, compression or retention;
\item additional batch allocations.
\end{itemize}
Network and WAN, Chilean hardware and staffing are not included. Why the 2026
forecast costs more than the construction-era one, although demand is mostly
lower, is explained in \secref{sec:upwhy}.

% GENERATED by generate_model.py from lsst/rubin-sizing-model @ v2026.10.0. Do not edit by hand.
% Regenerate: git clone --branch v2026.10.0 https://github.com/lsst/rubin-sizing-model ; then, inside the clone, uv run generate_model.py --emit-tex <this directory>

\small \begin{longtable}{|p{0.55\textwidth}|r|l|}
\caption{Unit costs and hardware lifetimes used by the cost forecast. \label{tab:smCostInputs}}\\
\hline
\textbf{Metric} & \textbf{Value} & \textbf{Unit} \\
\hline\endfirsthead
\hline
\textbf{Metric} & \textbf{Value} & \textbf{Unit} \\
\hline\endhead
Torino batch node (768 GB RAM) & \$89,355 & \$/node \\ \hline
K8s / Qserv-class node (EPYC 7543P) & \$30,866 & \$/node \\ \hline
Ceph usable disk & 235.05 & \$/TB \\ \hline
Tape / archive & 8.625 & \$/TB \\ \hline
Hosting per node per year & \$566 & \$/node-yr \\ \hline
Overhead on hardware & 14\% & \% \\ \hline
Overhead, fixed annual & \$150,000 & \$/yr \\ \hline
Annual price change & 0\% & \%/yr \\ \hline
Compute hardware lifetime & 5 & years \\ \hline
Storage hardware lifetime & 5 & years \\ \hline
K8s annual growth & 5\% & fraction \\ \hline
\end{longtable} \normalsize


\begin{landscape}
% GENERATED by generate_model.py from lsst/rubin-sizing-model @ v2026.10.0. Do not edit by hand.
% Regenerate: git clone --branch v2026.10.0 https://github.com/lsst/rubin-sizing-model ; then, inside the clone, uv run generate_model.py --emit-tex <this directory>

\tiny \begin{longtable}{|p{0.13\textwidth}|l|r|r|r|r|r|r|r|r|r|r|r|}
\caption{USDF hardware purchases and spend by operations year with all processing at the USDF, priced with USDF quotes. The first column uses actual current-year requests. Batch nodes serve DRP only; Prompt Processing is in the Kubernetes lines. \label{tab:smCost}}\\
\hline
\textbf{Metric} & \textbf{Unit} & \textbf{Year 1 (DP2, actual)} & \textbf{Year 2} & \textbf{Year 3} & \textbf{Year 4} & \textbf{Year 5} & \textbf{Year 6} & \textbf{Year 7} & \textbf{Year 8} & \textbf{Year 9} & \textbf{Year 10} & \textbf{10-yr Total} \\
 &  & FY2026 & FY2027 & FY2028 & FY2029 & FY2030 & FY2031 & FY2032 & FY2033 & FY2034 & FY2035 &  \\
\hline\endfirsthead
\hline
\textbf{Metric} & \textbf{Unit} & \textbf{Year 1 (DP2, actual)} & \textbf{Year 2} & \textbf{Year 3} & \textbf{Year 4} & \textbf{Year 5} & \textbf{Year 6} & \textbf{Year 7} & \textbf{Year 8} & \textbf{Year 9} & \textbf{Year 10} & \textbf{10-yr Total} \\
 &  & FY2026 & FY2027 & FY2028 & FY2029 & FY2030 & FY2031 & FY2032 & FY2033 & FY2034 & FY2035 &  \\
\hline\endhead
Batch nodes purchased & nodes & 0 & 0 & 72 & 86 & 77 & 42 & 42 & 114 & 128 & 119 & 680 \\ \hline
Batch compute cost & \$ & \$0 & \$0 & \$6,433,560 & \$7,684,530 & \$6,880,335 & \$3,752,910 & \$3,752,910 & \$10,186,470 & \$11,437,440 & \$10,633,245 & \$60,761,400 \\ \hline
Disk purchased this FY (new capacity) & TB & 10,082 & 0 & 19,231 & 31,081 & 31,081 & 31,081 & 31,081 & 31,081 & 31,081 & 31,081 & 246,883 \\ \hline
Disk re-buy (end of storage lifetime) & TB & 0 & 0 & 0 & 0 & 0 & 10,082 & 0 & 19,231 & 31,081 & 31,081 & 91,476 \\ \hline
Disk cost & \$ & \$2,369,794 & \$0 & \$4,520,240 & \$7,305,749 & \$7,305,749 & \$9,675,544 & \$7,305,749 & \$11,825,989 & \$14,611,499 & \$14,611,499 & \$79,531,812 \\ \hline
Tape / archive added (from All at USDF) & TB & 5,000 & 0 & 0 & 0 & 0 & 0 & 0 & 788 & 15,389 & 17,173 & 38,350 \\ \hline
Tape cost & \$ & \$43,125 & \$0 & \$0 & \$0 & \$0 & \$0 & \$0 & \$6,798 & \$132,731 & \$148,118 & \$330,772 \\ \hline
K8s nodes purchased (growth + cohort refresh) & nodes & 24 & 20 & 36 & 34 & 4 & 28 & 24 & 40 & 38 & 8 & 256 \\ \hline
K8s cost & \$ & \$740,784 & \$617,320 & \$1,111,176 & \$1,049,444 & \$123,464 & \$864,248 & \$740,784 & \$1,234,640 & \$1,172,908 & \$246,928 & \$7,901,696 \\ \hline
Qserv nodes purchased & nodes & 0 & 11 & 78 & 169 & 98 & 101 & 11 & 78 & 169 & 102 & 817 \\ \hline
Qserv cost & \$ & \$0 & \$339,526 & \$2,407,548 & \$5,216,354 & \$3,024,868 & \$3,117,466 & \$339,526 & \$2,407,548 & \$5,216,354 & \$3,148,332 & \$25,217,522 \\ \hline
Hardware subtotal & \$ & \$3,153,703 & \$956,846 & \$14,472,524 & \$21,256,077 & \$17,334,416 & \$17,410,168 & \$12,138,969 & \$25,661,445 & \$32,570,931 & \$28,788,121 & \$173,743,202 \\ \hline
Hosting & \$ & \$203,194 & \$220,740 & \$258,662 & \$362,806 & \$444,310 & \$541,096 & \$578,452 & \$624,864 & \$670,144 & \$700,708 & \$4,604,976 \\ \hline
Overhead & \$ & \$591,518 & \$283,958 & \$2,176,153 & \$3,125,851 & \$2,576,818 & \$2,587,423 & \$1,849,456 & \$3,742,602 & \$4,709,930 & \$4,180,337 & \$25,824,048 \\ \hline
TOTAL ANNUAL SPEND & \$M & \$3.95M & \$1.46M & \$16.91M & \$24.74M & \$20.36M & \$20.54M & \$14.57M & \$30.03M & \$37.95M & \$33.67M & \$204.17M \\ \hline
Cumulative spend & \$M & \$3.95M & \$5.41M & \$22.32M & \$47.06M & \$67.42M & \$87.96M & \$102.52M & \$132.55M & \$170.50M & \$204.17M &  \\ \hline
\end{longtable} \normalsize

% GENERATED by generate_model.py from lsst/rubin-sizing-model @ v2026.09.2. Do not edit by hand.
% Regenerate: git clone --branch v2026.09.2 https://github.com/lsst/rubin-sizing-model ; then, inside the clone, uv run generate_model.py --emit-tex <this directory>

\tiny \begin{longtable}{|p{0.13\textwidth}|l|r|r|r|r|r|r|r|r|r|r|r|}
\caption{USDF hardware purchases and spend with DRP processing shared at the Facility Split shares. Only the USDF is priced; partner facilities buy their own hardware. \label{tab:smCostSplit}}\\
\hline
\textbf{Metric} & \textbf{Unit} & \textbf{DP2 (actual)} & \textbf{DR1} & \textbf{DR2} & \textbf{DR3} & \textbf{DR4} & \textbf{DR5} & \textbf{DR6} & \textbf{DR7} & \textbf{DR8} & \textbf{DR9} & \textbf{10-yr Total} \\
 &  & FY2026 & FY2027 & FY2028 & FY2029 & FY2030 & FY2031 & FY2032 & FY2033 & FY2034 & FY2035 &  \\
\hline\endfirsthead
\hline
\textbf{Metric} & \textbf{Unit} & \textbf{DP2 (actual)} & \textbf{DR1} & \textbf{DR2} & \textbf{DR3} & \textbf{DR4} & \textbf{DR5} & \textbf{DR6} & \textbf{DR7} & \textbf{DR8} & \textbf{DR9} & \textbf{10-yr Total} \\
 &  & FY2026 & FY2027 & FY2028 & FY2029 & FY2030 & FY2031 & FY2032 & FY2033 & FY2034 & FY2035 &  \\
\hline\endhead
USDF peak concurrent cores (USDF share, Milano-equivalent) & cores & 794 & 2,016 & 4,032 & 10,885 & 14,514 & 18,142 & 21,771 & 25,399 & 29,028 & 32,656 &  \\ \hline
Batch nodes purchased & nodes & 0 & 0 & 72 & 50 & 35 & 0 & 1 & 87 & 65 & 51 & 361 \\ \hline
Batch compute cost & \$ & \$0 & \$0 & \$6,433,560 & \$4,467,750 & \$3,127,425 & \$0 & \$89,355 & \$7,773,885 & \$5,808,075 & \$4,557,105 & \$32,257,155 \\ \hline
On-floor disk need at USDF (less live intermediates held at partners) & TB & 6,733 & 16,170 & 30,909 & 56,310 & 76,030 & 95,750 & 115,470 & 135,190 & 154,910 & 174,630 &  \\ \hline
Disk purchased this FY (new capacity) & TB & 10,082 & 0 & 0 & 16,228 & 19,720 & 19,720 & 19,720 & 19,720 & 19,720 & 19,720 & 144,630 \\ \hline
Disk re-buy (end of storage lifetime) & TB & 0 & 0 & 0 & 0 & 0 & 10,082 & 0 & 0 & 16,228 & 19,720 & 46,030 \\ \hline
Disk cost & \$ & \$2,369,794 & \$0 & \$0 & \$3,814,443 & \$4,635,234 & \$7,005,028 & \$4,635,234 & \$4,635,234 & \$8,449,677 & \$9,270,468 & \$44,815,112 \\ \hline
Tape cost & \$ & \$43,125 & \$0 & \$0 & \$0 & \$0 & \$0 & \$0 & \$6,798 & \$132,731 & \$148,118 & \$330,772 \\ \hline
K8s cost & \$ & \$740,784 & \$617,320 & \$1,111,176 & \$1,049,444 & \$123,464 & \$864,248 & \$740,784 & \$1,234,640 & \$1,172,908 & \$246,928 & \$7,901,696 \\ \hline
Qserv cost & \$ & \$0 & \$339,526 & \$2,407,548 & \$5,216,354 & \$3,024,868 & \$3,117,466 & \$339,526 & \$2,407,548 & \$5,216,354 & \$3,148,332 & \$25,217,522 \\ \hline
Hardware subtotal & \$ & \$3,153,703 & \$956,846 & \$9,952,284 & \$14,547,991 & \$10,910,991 & \$10,986,742 & \$5,804,899 & \$16,058,105 & \$20,779,745 & \$17,370,951 & \$110,522,256 \\ \hline
Hosting & \$ & \$203,194 & \$220,740 & \$258,662 & \$342,430 & \$400,162 & \$473,176 & \$487,326 & \$518,456 & \$548,454 & \$564,302 & \$4,016,902 \\ \hline
Overhead & \$ & \$591,518 & \$283,958 & \$1,543,320 & \$2,186,719 & \$1,677,539 & \$1,688,144 & \$962,686 & \$2,398,135 & \$3,059,164 & \$2,581,933 & \$16,973,116 \\ \hline
TOTAL ANNUAL SPEND & \$M & \$3.95M & \$1.46M & \$11.75M & \$17.08M & \$12.99M & \$13.15M & \$7.25M & \$18.97M & \$24.39M & \$20.52M & \$131.51M \\ \hline
Cumulative spend & \$M & \$3.95M & \$5.41M & \$17.16M & \$34.24M & \$47.23M & \$60.38M & \$67.63M & \$86.61M & \$111.00M & \$131.51M &  \\ \hline
\end{longtable} \normalsize

\end{landscape}
